% Chapter-local macros (from papers/paper-lean-mlir/paper.tex preamble).
\newcommand{\messa}[1]{\ensuremath{\texttt{LeanMLIR}(#1)}}

\chapter{Verifying Peephole Rewriting in SSA Compiler IRs}
\label{ch:lean-mlir}

\paragraph*{Attribution.} This chapter is based on \emph{Verifying Peephole
Rewriting in SSA Compiler IRs}~\citep{leanmlir}, published at ITP 2024, joint
work with Alex Keizer, Chris Hughes, Andrés Goens, and Tobias Grosser.

\lettrine{S}{tatic} single assignment (SSA)~\citep{rastello2022ssa}
intermediate representations are at the core of modern compilers, thanks to
the benefits their immediate encoding of use-def relationships brings to
compiler analyses and transformations. \emph{Peephole
optimizations}~\citep{mckeeman1965peephole}, which replace instruction
sequences of bounded length with semantically equivalent optimized ones, are
so common that 10\% of all IR-transforming code in
LLVM~\citep{lattner2004llvm} belongs to its InstCombine peephole
optimizer---more than the size of LLVM's loop optimizer. The most referenced
tool in LLVM's commit log is the Alive peephole
verifier~\citep{lopes2021alive2}, which brought automatic SMT-based
verification into the day-to-day of the LLVM community.

There is, at the same time, an increasing need for \emph{domain-specific}
reasoning in modern compilers, fuelling tailored SSA-based IRs as popularized
by the MLIR compiler framework~\citep{lattner2020mlir}. Interactive theorem
provers provide strong guarantees for end-to-end compiler verification, as in
CompCert~\citep{leroy2016compcert}, but modern compilers and their IRs evolve
at a rate that makes proof engineering alongside them prohibitively expensive.
In this chapter we combine the convenience of automation with the versatility
of ITPs for verifying peephole rewrites across domain-specific IRs. We
formalize a core calculus for SSA-based IRs that is generic over the IR and
covers \emph{regions}---the nested scoping used by many domain-specific IRs in
the MLIR ecosystem---and mechanize it in Lean with a user-friendly frontend
that translates MLIR syntax into our calculus.

Concretely, this chapter contributes:
\begin{itemize}
  \item a formalization of SSA with regions, parametrized over a user-defined
  IR $X$, mechanized in our framework \messa{X}, exploiting
  denotational-style value semantics for optimizing along the SSA use-def
  chain of an MLIR-style IR;
  \item a verified peephole rewriter, with a proof that lifting a peephole
  rewrite to a rewrite on the entire program preserves semantics, together
  with verified dead code elimination and common subexpression elimination,
  and proof automation that eliminates the abstraction overhead of the SSA
  calculus;
  \item an extension of pure optimizations to contexts with side effects; and
  \item syntax, semantics, and local rewrites for three MLIR-based IRs:
  arithmetic over bitvectors, structured control flow, and fully homomorphic
  encryption.
\end{itemize}

\section{Motivation: Verifying Optimizations for High-Level IRs}
\label{sec:mlir:motivation}

\sid{Port from papers/paper-lean-mlir/paper.tex \S 2 (lines 505--620).}

\section{\texorpdfstring{\messa{X}}{LeanMLIR(X)}: A Framework for Intrinsically Well-Typed SSA}
\label{sec:mlir:framework}

\sid{Port from papers/paper-lean-mlir/paper.tex \S 3 (lines 621--760).}

\section{Reasoning About \texorpdfstring{\messa{X}}{LeanMLIR(X)}}
\label{sec:mlir:reasoning}

\sid{Port from papers/paper-lean-mlir/paper.tex \S 4 (lines 761--863),
including the verified peephole rewriter, DCE, CSE, and proof automation.}

\section{Pure Rewriting in a Side-Effectful World}
\label{sec:mlir:effects}

\sid{Port from papers/paper-lean-mlir/paper.tex \S 5 (lines 864--890).}

\section{Case Studies}
\label{sec:mlir:casestudies}

\sid{Port from papers/paper-lean-mlir/paper.tex \S 6 (lines 891--1146):
bitvector rewrites from LLVM, structured control flow, fully homomorphic
encryption. Figures live in chapters/lean-mlir/plots/ and
chapters/lean-mlir/images/; the Ott-generated grammar is at
chapters/lean-mlir/ott/lc.nowrap.tex.}

\section{Related Work}
\label{sec:mlir:related}

\sid{Port from papers/paper-lean-mlir/paper.tex \S 7 (lines 1147--1181).}

\section{Conclusions}
\label{sec:mlir:conclusion}

\sid{Port from papers/paper-lean-mlir/paper.tex \S 8 (lines 1182--), and add
a hand-off paragraph: verified rewriting exposes bitvector proof obligations,
motivating \cref{ch:single-width-bv}.}
